\ssr{ВВЕДЕНИЕ}

Целью работы является выполнение сравнительного анализа ресурсной эффективности алгоритмов поиска элементов в массиве.

Для достижения поставленной цели были определены следующие задачи:
\begin{itemize}
	\item описать решаемую согласно варианту задачу;
	\item разработать рекурсивный и нерекурсивные алгоритмы решения задачи, согласно варианту;
	\item описать разработанные алгоритмы схемами алгоритмов;
	\item описать требования к программному обеспечению;
	\item описать средства реализации алгоритмов;
	\item реализовать разработанные алгоритмы;
	\item выполнить тестирование реализации алгоритмов;
	\item выполнить теоретическую оценку затрачиваемой реализацией каждого алгоритма памяти (для рекурсивного алгоритма~--- на материале анализа высоты дерева рекурсивных вызовов и оценки затрачиваемой на один вызов функции памяти);
	\item оценить временную эффективность алгоритмов или их реализаций в лучшем и худшем случае в терминах числа сравнений искомой величины с другими величинами;
	\item представить результаты сравнительного анализа в формате таблицы;
	\item сделать выводы из сравнительного анализа реализации разработанных алгоритмов решения одной и той же задачи, заданной вариантом, по критериям временной и ёмкостной эффективности, а также дополнительных требований, налагаемых алгоритмом.
\end{itemize}